\section{Methods}
+ scGPT erklären wie es funktioniert/trainiert
+ scGPT wieso? hypothese staten
\subsection{Data}

All experiments use the highest-confidence intersection of SCP542 \cite{scp542} and CTRPv2 \cite{ctrpv2,ctrp_rees}. This intersection comprises \NLines{} cell lines with
post-quality-control measurements (190 roster name matches, of which 10 are unscreened), \NDrugs{}
compounds and \NCells{} cells, with effectively complete viability coverage inside the overlap. Genes are
reduced to the \NHvg{} most highly variable; scGPT embeds the \NVocab{} of these that lie in its
vocabulary, while PCA is computed on all \NHvg{}. The gene-set size is not critical: an earlier sweep from
1{,}000 genes to the full transcriptome showed no advantage for any particular count
(\code{notebooks/07}), so the \NHvg-gene selection is used throughout.
\revision{\emph{Correction, 05.08.2026.} That sweep was computed on the legacy \code{mean\_pv} response
score rather than the \code{auc} target used throughout this report, so its numbers are superseded; it
now lives in \code{notebooks/data\_and\_harmonization/verify\_variants.ipynb}, \S9, and the reference to
\code{notebooks/07} above is stale. Its \emph{all-genes} arm is also not the full transcriptome as far as
scGPT is concerned: the model reads at most 1{,}200 genes per cell, drawn at random from that cell's
non-zero genes, so what the sweep compared was \NHvg{} \emph{selected} genes against a random subsample
of broadly similar size. The claim that gene-set size does not matter is therefore untested on the
current target.}

\begin{revisionblock}
\textbf{Expression preprocessing (revised 05.08.2026).} SCP542 is distributed as counts per million, so
library-size normalization has already been applied --- once, to the full gene matrix, which is where it
belongs. The pipeline therefore does not normalize a second time. Expression is log-transformed as
$\log(1+x)$; the \NHvg{} most variable genes are selected on that matrix using the dispersion-based
Seurat criterion; each gene is then standardized to zero mean and unit variance across cells, with
$z$-scores clipped at 10, before the principal components are taken
(\code{scripts/preprocessing/scp542\_conversion.py}, \code{scripts/preprocessing/add\_pca.py}).

Until this revision the PCA branch did two things differently. It applied a second library-size
normalization \emph{after} gene selection, dividing each cell by the total of its retained genes alone,
which rescales a cell by the share of its expression that happens to fall inside the selected set and
scales the same cell differently in each gene-set variant. And it applied no per-gene standardization,
so the leading components tracked absolute expression level rather than variation between cells. Both
are corrected in the code; \textbf{every PCA number in this report was produced before the correction}
and will change when the representation is recomputed.

scGPT receives the same CPM matrix with neither step applied. It replaces them with a per-cell binning
of expression into 51 quantile bins, computed from that cell's own non-zero values --- a rank transform,
and therefore invariant to any monotone per-cell rescaling. Counts, CPM and $\log(1+x)$ inputs give
identical bins, so the CPM input does not disadvantage the model \cite{scgpt}.
\end{revisionblock}

A property of the data governs every error estimate in this report. The label is defined per cell line and
drug, not per cell, so all cells of a line share the same broadcast bulk value and act as
pseudo-replicates. The effective sample size is therefore the number of cell lines, approximately
\NIndep{} after the held-out splits are removed, not the \NCells{} cells. Treating cells as independent
would understate the uncertainty by roughly the square root of the number of cells per line.

\subsection{Representation and model}

The comparison between representations is controlled so that only the content of the embedding differs.
Both the PCA baseline (\code{X\_pca}) and the scGPT embedding (\code{X\_scGPT}) are 512-dimensional and are
fed into the same multi-head perceptron, with hidden dimensions 128 and 64, layer normalization, GELU
activations and dropout. The network has one output head per drug and is trained with a masked
mean-squared-error loss that is evaluated only on observed cell-drug pairs. The model and training code are
defined in \code{scripts/model/OncoMLP.py} and \code{scripts/training/train\_multitask.py}.

\begin{revisionblock}
\textbf{Splits, and where each representation is fitted (added 05.08.2026).} Train, validation and test
are a 70/15/15 partition of cell lines, so every cell of a line falls on the same side and no line is
split across parts. The assignment is now \emph{frozen to a versioned file}
(\code{splits/split\_ctrp.csv}) instead of being redrawn from the data on each run. It has to be:
eligibility depends on which compounds survive the label filters, so a change to the drug panel would
otherwise move lines between train, validation and test unannounced, and results from either side of
that change would look comparable while being scored on different held-out lines. Redrawing requires an
explicit flag, and a line that the frozen file does not cover is an error rather than a fresh draw.

The principal components used as model input are fitted on the training lines of that split alone ---
per-gene standardization and the rotation both --- and then applied to the held-out lines
(\code{X\_pca\_train\_ctrp}). A separate all-cells decomposition (\code{X\_pca}) is kept for description
only, such as the latent-space figures, where holding cells out would be wrong. The cross-validation
results are the exception: its folds are drawn at training time rather than stored, so no precomputed
train-only projection exists for them and they still use the all-cells decomposition. scGPT needs no
such distinction --- its weights are pretrained and frozen, and its per-cell binning uses only the cell
being embedded, so nothing about it is fitted on this data.
\end{revisionblock}

\subsection{Response target and per-drug standardization}
\label{sec:target}

The choice of training target requires care. CTRPv2 provides a dose-averaged percent-viability score,
denoted \code{mean\_pv}, and an integrated area under the fitted dose-response curve, \code{auc}. The
initial experiments used \code{mean\_pv}; the final target is the curve-fit AUC standardized per drug
across cell lines, \code{auc\_z}, defined for cell line $i$ and drug $j$ as
$\mathrm{auc\_z}_{ij} = (\mathrm{auc}_{ij} - \mu_j)/\sigma_j$, where $\mu_j$ and $\sigma_j$ are the mean and
standard deviation of that drug's AUC over the cell lines.

Two considerations motivate the standardization. First, the per-drug response variance is highly
heterogeneous across the \NDrugs{} compounds, with the per-drug standard deviation ranging from 0.034 to
0.302, a factor of roughly 80 in squared error. In a shared multi-task objective with one masked
squared-error term per drug, the contribution of a drug scales with its variance, so a small number of
high-variance compounds dominate the shared representation irrespective of whether that variance carries
any learnable signal; the three drugs with the largest spread (\code{ifosfamide}, \code{fqi-2} and
\code{bafilomycin a1}) in fact fail to separate the cell lines at all. Standardizing each drug to unit
variance is algebraically equivalent to weighting each head by $1/\sigma_j^2$, which removes this
imbalance and is the regression analogue of loss reweighting under class imbalance \cite{kendall}. Second,
per-drug standardization changes the question from absolute potency, which is a property of the drug, to
the relative sensitivity of a cell line for a given drug, which is the quantity of interest for
personalized prediction. Because the Spearman correlation within a drug is invariant to an affine per-drug
transform, the standardization does not alter any single-drug evaluation metric; its only effect is on the
multi-task coupling. The target is implemented in \code{scripts/preprocessing/ctrp\_to\_h5ad.py}.

\subsection{Evaluation}

Cell lines are evaluated out of fold. A 5-fold group $k$-fold split, grouped by cell line, is applied to
the 153 training and validation lines, so that no cell of a test line is seen during training. Each line
is predicted exactly once, the per-cell predictions of a line are averaged to the line level, and the
Spearman correlation between predicted and true values across the roughly \NIndep{} lines is computed per
drug and averaged over drugs. This protocol replaces an earlier fixed 27-line validation split, on which
the standard error of a correlation is approximately 0.2 and effects below 0.4 are not measurable. The
Pearson correlation tracks the Spearman value within 0.02 throughout.

For external comparison, the model is additionally evaluated with the DrEval package \cite{dreval}
(version 1.5.1), using its leave-cell-line-out splits, its baselines and its evaluation routine without
reimplementation. As recommended by the framework, metrics are reported after the mean drug and cell-line
effects have been removed from both the true and predicted responses, and they are pooled over all
cell-line-drug pairs rather than averaged per drug, so they are not directly comparable to the per-drug
correlations above.
